% partes/parte3/intro.tex
\part{Modelos de comportamiento del consumidor digital}

\chapter*{Introducción a la Parte III}
\addcontentsline{toc}{chapter}{Introducción a la Parte III}

La Parte II describió cómo decide una persona. La Parte III pregunta
\emph{con qué teorías} se explica esa decisión. El programa oficial la
llama \enquote{Modelos de comportamiento del consumidor digital}, y su
competencia es distinguir los modelos de estudio a partir de los tipos de
comportamiento y del análisis de los modelos de decisión de compra.

El Capítulo~9 abre la Parte con dos bloques. El primero es una tipología de
cuatro formas de comprar (compleja, con reducción de disonancia, habitual y
variada) que se distinguen por el involucramiento de la persona y por cuánto
difieren las marcas a sus ojos. El segundo recorre nueve modelos de estudio, desde
el económico de Marshall hasta la netnografía, y los aplica al mismo caso para
que se note qué ve cada uno y qué deja sin ver. Cada modelo se presenta con su
supuesto, un ejemplo digital y su límite. Dos de ellos (el psicoanalítico y el
motivacional) exigen corregir lo que los manuales les atribuyen, y uno (la
netnografía) es un método, no una teoría. El tema de la base tecnológica de la
unidad (programa 3.3) se tratará en los capítulos siguientes de la Parte.

La advertencia metodológica de la Parte II sigue vigente: los modelos son
representaciones útiles y estilizadas, la mayoría nació fuera de la
mercadotecnia y ninguno anticipó el comercio electrónico. Se presentan como
lo que son.
